\section{Segundo ensaio: custo-benefício e escolha do instrumento}
\label{sec:cb}

O primeiro ensaio mede o que a lei fez. O segundo faz duas perguntas sobre
esses números: se o benefício em saúde supera o custo em produção, e se outro
instrumento teria feito melhor.

\subsection{Custo-benefício}

Os benefícios são o efeito sobre a saúde ao nascer, valorado pela literatura de
retornos ao peso ao nascer e pelo custo hospitalar neonatal, e a variação nas
intoxicações, valorada pelo método do custo da doença de
\citet{soares2009socialcost, soares2012impactos}. Os valores publicados por
esses autores são do Paraná em 1998 e 1999 e servem de método; os preços serão
reestimados com as internações do Ceará e do Rio Grande do Norte. Os custos são
a variação de área, produção e emprego estimada no primeiro ensaio e a
diferença de custo entre as tecnologias de aplicação.

Dois cuidados. O benefício calculado é um piso, porque não inclui efeitos
crônicos nem o dano ambiental. E a substituição entra dos dois lados da conta:
se a troca do avião pela aplicação manual elevar a intoxicação de
trabalhadores, esse aumento é custo da política.

Sobre o custo há, hoje, duas afirmações em conflito e nenhuma estimativa com
contrafactual. Produtores previram perda de metade da produção. Uma nota
técnica mostra que a produção e a exportação de banana do Ceará subiram depois
de 2018 \citep{cavalcante2022nota}; a série do IBGE confirma o fato, e mostra
também que 2018 é o fundo de uma seca de seis anos e que o Rio Grande do Norte
permaneceu estável. O contraste entre os dois lados da Chapada será a primeira
estimativa com grupo de comparação.

O resultado será apresentado em cenários, com a pergunta invertida: que efeito
sobre o peso ao nascer seria necessário para a lei se pagar, dado o custo
estimado, e esse efeito está dentro ou fora do intervalo de confiança?

\subsection{O instrumento}

A teoria organiza a comparação. Sob incerteza, a regulação por quantidade supera
a por preço quando o dano marginal cresce mais depressa que o custo marginal de
evitá-lo \citep{weitzman1974}. A deriva é poluição difusa, sem medidor; uma
taxa sobre ela não é implementável, e uma taxa sobre o voo não distingue o
bananal ao lado da escola do bananal isolado. A regulação de agrotóxicos tem,
além disso, efeitos distributivos que dependem de quem consegue substituir
\citep{lichtenberg1988}.

Quatro fatos do caso brasileiro informam a discussão. A zona de exclusão já
existia na norma federal e era descumprida. O instrumento de preço opera com o
sinal trocado: a desoneração fiscal dos agrotóxicos somou R\$~9,86 bilhões em
2017 \citep{cunha2020incentivos}. A produtividade marginal do agrotóxico é
baixa em quase todo o país e mais alta no Nordeste \citep{rodrigues2022tese}, de
modo que um instrumento nacional de preço seria mal calibrado para a região. E
o legislador cearense já se moveu, em 2024, da proibição para a regulação do
método, ao admitir drones.

O ensaio registra de antemão o que conclui em cada resultado do primeiro
(Tabela~\ref{tab:instrumento}), para que a interpretação não se ajuste ao
número.

\begin{table}[ht]
\centering\small
\caption{Leitura do instrumento conforme o resultado do primeiro ensaio}
\label{tab:instrumento}
\begin{tabularx}{\textwidth}{@{}LL@{}}
\toprule
Se o primeiro ensaio encontrar & Leitura \\
\midrule
Melhora ao nascer, sem perda de produção & A proibição foi barata e eficaz; uma taxa teria mantido voos junto às mesmas comunidades \\
Melhora ao nascer, com perda de produção & O custo-benefício decide; a zona de exclusão com fiscalização efetiva é o candidato a dominar \\
Nenhum efeito detectável ao nascer; intoxicações sobem & A lei trocou um risco difuso por um ocupacional; o problema era a molécula e a prática \\
Nenhum efeito detectável, e o avião já não voava & A lei consolidou um estado que a seca havia produzido; seu valor é o de impedir a retomada, e o desenho não o enxerga \\
\bottomrule
\end{tabularx}
\end{table}
